% ============================================================
% Question Bank: ch08-03 Comparing Two Populations Visually (Modified — FL/IN/MN/VA)
% Topic Title: Comparing Two Populations Visually (Multiple Display Emphasis)
% CCSS: 7.SP.B.3 (FL/IN/MN/VA: multiple display emphasis)
% Questions: 30 (18 MC + 7 SA + 5 Graphical)
% ============================================================

% --- Multiple Choice Questions (18) ---

\begin{practiceQuestion}{m-8-3-q01}{mc}
\begin{questionText}
Which of the following data displays keeps every original data value while also showing the shape of the distribution?
\end{questionText}
\choiceA{Histogram}
\choiceB{Stem-and-leaf plot}
\choiceC{Box plot}
\choiceD{Circle graph}
\correctAnswer{B}
\explanation{A stem-and-leaf plot preserves every data value (read from stem and leaf) and shows the overall shape.}
\end{practiceQuestion}

\begin{practiceQuestion}{m-8-3-q02}{mc}
\begin{questionText}
When is a box plot a better choice than a dot plot for comparing two groups?
\end{questionText}
\choiceA{When data sets are very small (fewer than $10$ values)}
\choiceB{When you want to see every individual data point}
\choiceC{When data sets are large and you want to compare medians and variability quickly}
\choiceD{When the data is categorical}
\correctAnswer{C}
\explanation{Box plots summarize large data sets with five key values (min, $Q_1$, median, $Q_3$, max), making them efficient for comparing groups.}
\end{practiceQuestion}

\begin{practiceQuestion}{m-8-3-q03}{mc}
\begin{questionText}
A circle graph shows Class A spends $40\%$ of free time on sports and Class B spends $25\%$. What can you conclude?
\end{questionText}
\choiceA{More students in Class A play sports}
\choiceB{Class A devotes a larger proportion of free time to sports}
\choiceC{Class B likes sports less}
\choiceD{The classes have the same number of students}
\correctAnswer{B}
\explanation{Circle graphs compare proportions (parts of a whole), not actual counts. Class A devotes a larger share to sports.}
\end{practiceQuestion}

\begin{practiceQuestion}{m-8-3-q04}{mc}
\begin{questionText}
Two groups have overlapping box plots. The medians differ by $3$ and both IQRs are $10$. What does this tell you?
\end{questionText}
\choiceA{The groups are completely different}
\choiceB{The difference is meaningful because $3 > 0$}
\choiceC{The difference is small relative to the variability — there is much overlap}
\choiceD{Box plots cannot be used to compare groups}
\correctAnswer{C}
\explanation{A difference of $3$ with IQR of $10$ means the overlap is large. The difference is not very meaningful.}
\end{practiceQuestion}

\begin{practiceQuestion}{m-8-3-q05}{mc}
\begin{questionText}
Which display is best for comparing the proportions (parts of a whole) for two groups?
\end{questionText}
\choiceA{Box plot}
\choiceB{Circle graph}
\choiceC{Dot plot}
\choiceD{Stem-and-leaf plot}
\correctAnswer{B}
\explanation{Circle graphs show how a whole is divided into parts, making them ideal for comparing proportions between groups.}
\end{practiceQuestion}

\begin{practiceQuestion}{m-8-3-q06}{mc}
\begin{questionText}
A frequency table for Group A shows: $60$--$69$: $2$, $70$--$79$: $8$, $80$--$89$: $10$, $90$--$99$: $5$. What is the total number of data values?
\end{questionText}
\choiceA{$4$}
\choiceB{$20$}
\choiceC{$25$}
\choiceD{$99$}
\correctAnswer{C}
\explanation{$2 + 8 + 10 + 5 = 25$.}
\end{practiceQuestion}

\begin{practiceQuestion}{m-8-3-q07}{mc}
\begin{questionText}
A back-to-back stem-and-leaf plot is useful because it:
\end{questionText}
\choiceA{Hides individual data values}
\choiceB{Shows two data sets side by side using the same stems}
\choiceC{Only works for data sets of exactly the same size}
\choiceD{Is the only way to compare two groups}
\correctAnswer{B}
\explanation{A back-to-back plot shares stems, with one group's leaves on the left and the other's on the right, enabling direct comparison.}
\end{practiceQuestion}

\begin{practiceQuestion}{m-8-3-q08}{mc}
\begin{questionText}
Group X: mean $70$, range $15$. Group Y: mean $75$, range $35$. Which statement is best?
\end{questionText}
\choiceA{Group Y scores higher and is more consistent}
\choiceB{Group Y scores higher but is less consistent}
\choiceC{Group X is both higher and more consistent}
\choiceD{The groups are the same}
\correctAnswer{B}
\explanation{Group Y has a higher mean ($75$ vs.\ $70$) but a much larger range ($35$ vs.\ $15$), meaning its data is more spread out—less consistent.}
\end{practiceQuestion}

\begin{practiceQuestion}{m-8-3-q09}{mc}
\begin{questionText}
A histogram shows that Group A's data is clustered around $80$--$90$ while Group B's data is spread evenly from $60$ to $100$. Both have a mean near $80$. Which group has a larger MAD?
\end{questionText}
\choiceA{Group A}
\choiceB{Group B}
\choiceC{They have the same MAD}
\choiceD{Cannot be determined from histograms}
\correctAnswer{B}
\explanation{Group B is spread evenly over a wider range, so its data points are farther from the mean on average, giving a larger MAD.}
\end{practiceQuestion}

\begin{practiceQuestion}{m-8-3-q10}{mc}
\begin{questionText}
When comparing two populations, why is it helpful to use more than one type of display?
\end{questionText}
\choiceA{It wastes time but is required}
\choiceB{Each display type reveals different aspects of the data}
\choiceC{Using more displays always changes the conclusion}
\choiceD{One display is always sufficient}
\correctAnswer{B}
\explanation{Dot plots show individual values, box plots show spread and center, histograms show shape, stem-and-leaf plots preserve values. Together they tell a more complete story.}
\end{practiceQuestion}

\begin{practiceQuestion}{m-8-3-q11}{mc}
\begin{questionText}
A frequency table shows categorical data for two classes: favorite sport. Which display would best compare these results?
\end{questionText}
\choiceA{A box plot}
\choiceB{A dot plot}
\choiceC{A double bar graph}
\choiceD{A number line}
\correctAnswer{C}
\explanation{A double bar graph compares frequencies of categorical data (like sports choices) for two groups side by side.}
\end{practiceQuestion}

\begin{practiceQuestion}{m-8-3-q12}{mc}
\begin{questionText}
Two dot plots for test scores show that Group A is centered around $85$ with most values between $80$ and $90$, while Group B is centered around $75$ with values between $60$ and $90$. Which group is more consistent?
\end{questionText}
\choiceA{Group B because it has a wider range of scores}
\choiceB{Group A because its data is more tightly clustered}
\choiceC{They are equally consistent}
\choiceD{Cannot tell from dot plots}
\correctAnswer{B}
\explanation{Group A's data clusters tightly ($80$--$90$), while Group B spreads from $60$ to $90$. Tighter clustering means more consistency.}
\end{practiceQuestion}

\begin{practiceQuestion}{m-8-3-q13}{mc}
\begin{questionText}
A stem-and-leaf plot shows exactly the same information as which other display?
\end{questionText}
\choiceA{Box plot}
\choiceB{Histogram}
\choiceC{An ordered list of all data values}
\choiceD{A circle graph}
\correctAnswer{C}
\explanation{A stem-and-leaf plot preserves every data value, just like a sorted list, while also showing the shape of the distribution.}
\end{practiceQuestion}

\begin{practiceQuestion}{m-8-3-q14}{mc}
\begin{questionText}
Two box plots overlap by half their boxes. The medians are $4$ apart and the IQRs are about $8$ each. Is the difference in centers meaningful?
\end{questionText}
\choiceA{Yes — any difference in medians is meaningful}
\choiceB{No — the difference ($4$) is only about half the IQR ($8$), showing much overlap}
\choiceC{Yes — because the IQR is $8$}
\choiceD{Cannot tell}
\correctAnswer{B}
\explanation{When the difference in centers is less than the IQR, there is substantial overlap and the difference is considered small.}
\end{practiceQuestion}

\begin{practiceQuestion}{m-8-3-q15}{mc}
\begin{questionText}
Class A: median $90$, IQR $5$. Class B: median $82$, IQR $12$. How do the classes compare?
\end{questionText}
\choiceA{Class A has a higher center and more variability}
\choiceB{Class B has a higher center and less variability}
\choiceC{Class A has a higher center and less variability}
\choiceD{The classes are the same}
\correctAnswer{C}
\explanation{Class A has a higher median ($90$ vs.\ $82$) and a smaller IQR ($5$ vs.\ $12$), meaning higher scores that are more tightly grouped.}
\end{practiceQuestion}

\begin{practiceQuestion}{m-8-3-q16}{mc}
\begin{questionText}
Which display type does NOT show the exact individual data values?
\end{questionText}
\choiceA{Dot plot}
\choiceB{Box plot}
\choiceC{Stem-and-leaf plot}
\choiceD{Frequency table (ungrouped)}
\correctAnswer{B}
\explanation{A box plot shows five summary statistics (min, $Q_1$, median, $Q_3$, max) but not individual data points.}
\end{practiceQuestion}

\begin{practiceQuestion}{m-8-3-q17}{mc}
\begin{questionText}
A student creates both a dot plot and a histogram for the same data set. The dot plot shows a cluster of values at $5$ and $6$, while the histogram's tallest bar is the $5$--$7$ interval. Are these consistent?
\end{questionText}
\choiceA{No — they contradict each other}
\choiceB{Yes — both show the most data near $5$--$7$}
\choiceC{No — histograms are never accurate}
\choiceD{They cannot be compared}
\correctAnswer{B}
\explanation{The dot plot clusters at $5$--$6$ and the histogram's peak at $5$--$7$ are consistent. The histogram groups data into intervals, slightly broadening the picture.}
\end{practiceQuestion}

\begin{practiceQuestion}{m-8-3-q18}{mc}
\begin{questionText}
Group A's scores range from $50$ to $95$; Group B's scores range from $70$ to $90$. Both groups have a median of $80$. Which group has more consistent performance?
\end{questionText}
\choiceA{Group A because its range is wider}
\choiceB{Group B because its range is narrower}
\choiceC{They are equally consistent because they have the same median}
\choiceD{Cannot be determined}
\correctAnswer{B}
\explanation{Group B's scores span only $20$ points ($70$--$90$) vs.\ Group A's $45$ points ($50$--$95$). Narrower range means more consistency.}
\end{practiceQuestion}

% --- Short Answer Questions (7) ---

\begin{practiceQuestion}{m-8-3-q19}{sa}
\begin{questionText}
Name three different types of displays that can be used to compare two data sets visually.
\end{questionText}
\correctAnswer{Dot plots, box plots, and stem-and-leaf plots (or histograms, frequency tables, circle graphs)}
\explanation{Multiple display types exist: dot plots, box plots, histograms, stem-and-leaf plots, circle graphs, and frequency tables each show data differently.}
\end{practiceQuestion}

\begin{practiceQuestion}{m-8-3-q20}{sa}
\begin{questionText}
Group P: mean $72$, MAD $6$. Group Q: mean $84$, MAD $6$. How many MADs apart are the means? Is the difference meaningful?
\end{questionText}
\correctAnswer{$2$ MADs; yes, the difference is meaningful}
\explanation{$\frac{84-72}{6} = 2$ MADs. Since this is $\geq 2$, the difference is considered meaningful with little overlap.}
\end{practiceQuestion}

\begin{practiceQuestion}{m-8-3-q21}{sa}
\begin{questionText}
Give one advantage of using a dot plot instead of a box plot when comparing two small groups (fewer than $15$ values each).
\end{questionText}
\correctAnswer{A dot plot shows every individual data point, so you can see the full distribution and any clusters or gaps}
\explanation{With small data sets, seeing every value is important. Box plots hide individual points, which can mask important features of small data sets.}
\end{practiceQuestion}

\begin{practiceQuestion}{m-8-3-q22}{sa}
\begin{questionText}
Class A: range $18$, IQR $8$. Class B: range $30$, IQR $8$. Both have the same IQR. Explain what the different ranges tell you.
\end{questionText}
\correctAnswer{Class B has more extreme values (outliers) that stretch the range, but the middle $50\%$ is equally spread in both classes}
\explanation{Same IQR means the central data is similarly spread. Class B's wider range means its extreme values are farther from the center.}
\end{practiceQuestion}

\begin{practiceQuestion}{m-8-3-q23}{sa}
\begin{questionText}
When would a circle graph be a useful tool for comparing two populations?
\end{questionText}
\correctAnswer{When comparing how each group divides a whole into proportional parts (e.g., comparing how two classes spend their time)}
\explanation{Circle graphs compare proportions. They are useful when you want to see what fraction of a whole each category represents for different groups.}
\end{practiceQuestion}

\begin{practiceQuestion}{m-8-3-q24}{sa}
\begin{questionText}
Data set A: $5, 7, 8, 9, 11$. Data set B: $3, 6, 8, 10, 13$. Find the median and range for each. Which set is more variable?
\end{questionText}
\correctAnswer{Both medians $= 8$. Range A $= 6$, Range B $= 10$. Set B is more variable.}
\explanation{Median A: $8$. Median B: $8$. Range A: $11 - 5 = 6$. Range B: $13 - 3 = 10$. Set B has a larger range, so it is more variable.}
\end{practiceQuestion}

\begin{practiceQuestion}{m-8-3-q25}{sa}
\begin{questionText}
A student says: ``I only need to look at the means to compare two groups.'' Is this correct? Explain.
\end{questionText}
\correctAnswer{No. You also need to look at variability (MAD, IQR, or range) to understand how spread out each group is.}
\explanation{Two groups can have the same mean but very different spreads. Comparing both center and variability gives a complete picture.}
\end{practiceQuestion}

% --- Graphical Questions (3 GMC + 2 GSA) ---

\begin{practiceQuestion}{m-8-3-q26}{gmc}
\begin{questionText}
The back-to-back stem-and-leaf plot compares two teams' scores.

\begin{center}
\begin{tabular}{r|c|l}
\textbf{Team A} & \textbf{Stem} & \textbf{Team B} \\
\hline
$9\;\;7\;\;5$ & $5$ & $2\;\;6\;\;8$ \\
$8\;\;4\;\;1\;\;0$ & $6$ & $0\;\;3\;\;5\;\;7$ \\
$6\;\;3$ & $7$ & $1\;\;4\;\;8\;\;9$ \\
$0$ & $8$ & $2\;\;5$
\end{tabular}

Key: $5 \mid 5$ means $55$
\end{center}

Which team has a higher median?
\end{questionText}
\choiceA{Team A}
\choiceB{Team B}
\choiceC{They have the same median}
\choiceD{Cannot be determined}
\correctAnswer{B}
\explanation{Team A ($10$ values): $55, 57, 59, 60, 61, 64, 68, 73, 76, 80$. Median $= \frac{61+64}{2} = 62.5$. Team B ($13$ values): $52, 56, 58, 60, 63, 65, 67, 71, 74, 78, 79, 82, 85$. Median $= 67$. Team B is higher.}
\end{practiceQuestion}

\begin{practiceQuestion}{m-8-3-q27}{gmc}
\begin{questionText}
The box plots compare daily high temperatures (°F) for City X and City Y over one month.

\begin{center}
\begin{tikzpicture}[scale=0.12]
  \draw[thick,->] (39,0) -- (91,0);
  \foreach \x in {40,45,...,90} {
    \draw (\x,0.5) -- (\x,-0.5);
    \node[below,font=\small] at (\x,-0.5) {$\x$};
  }
  % City X: min=50, Q1=58, med=65, Q3=72, max=82
  \draw[thick] (50,5) -- (82,5);
  \draw[thick,fill=funBlue!30] (58,3.5) rectangle (72,6.5);
  \draw[very thick,funBlue] (65,3.5) -- (65,6.5);
  \draw[thick] (50,3.5) -- (50,6.5);
  \draw[thick] (82,3.5) -- (82,6.5);
  \node[right,font=\small] at (83,5) {City X};
  % City Y: min=42, Q1=50, med=58, Q3=66, max=78
  \draw[thick] (42,12) -- (78,12);
  \draw[thick,fill=funOrange!30] (50,10.5) rectangle (66,13.5);
  \draw[very thick,funOrange] (58,10.5) -- (58,13.5);
  \draw[thick] (42,10.5) -- (42,13.5);
  \draw[thick] (78,10.5) -- (78,13.5);
  \node[right,font=\small] at (79,12) {City Y};
\end{tikzpicture}
\end{center}

Which city has warmer temperatures and which is more consistent?
\end{questionText}
\choiceA{City X is warmer; City Y is more consistent}
\choiceB{City X is warmer; both are equally consistent}
\choiceC{City X is warmer and more consistent}
\choiceD{City Y is warmer; City X is more consistent}
\correctAnswer{B}
\explanation{City X median $= 65°$, IQR $= 14$. City Y median $= 58°$, IQR $= 16$. City X is warmer. IQRs are similar ($14$ vs.\ $16$), so variability is about the same.}
\end{practiceQuestion}

\begin{practiceQuestion}{m-8-3-q28}{gmc}
\begin{questionText}
The frequency tables show quiz scores for two groups.

\begin{center}
\begin{tabular}{|c|c|c|}
\hline
\textbf{Score} & \textbf{Group A} & \textbf{Group B} \\
\hline
$60$--$69$ & $2$ & $5$ \\
\hline
$70$--$79$ & $6$ & $8$ \\
\hline
$80$--$89$ & $10$ & $4$ \\
\hline
$90$--$100$ & $7$ & $3$ \\
\hline
\end{tabular}
\end{center}

Which group has a higher proportion of scores $80$ or above?
\end{questionText}
\choiceA{Group A ($68\%$)}
\choiceB{Group B ($35\%$)}
\choiceC{Both are the same}
\choiceD{Group A ($50\%$)}
\correctAnswer{A}
\explanation{Group A total: $25$. Scores $\geq 80$: $10+7 = 17$. $\frac{17}{25} = 68\%$. Group B total: $20$. Scores $\geq 80$: $4+3 = 7$. $\frac{7}{20} = 35\%$. Group A is higher.}
\end{practiceQuestion}

\begin{practiceQuestion}{m-8-3-q29}{gsa}
\begin{questionText}
The dot plots show the number of hours worked per week by employees at two stores.

\begin{center}
\begin{tikzpicture}[scale=0.5]
  % Store A
  \node[left,font=\small\bfseries] at (14,2) {Store A};
  \draw[thick,->] (14.5,0) -- (26.5,0);
  \foreach \x in {15,16,...,26} {
    \draw (\x,0.1) -- (\x,-0.1);
    \node[below,font=\tiny] at (\x,0) {$\x$};
  }
  % Store A: 18(1),19(2),20(4),21(3),22(2),23(1) = 13
  \fill[funBlue] (18,0.35) circle (0.15);
  \fill[funBlue] (19,0.35) circle (0.15); \fill[funBlue] (19,0.7) circle (0.15);
  \fill[funBlue] (20,0.35) circle (0.15); \fill[funBlue] (20,0.7) circle (0.15); \fill[funBlue] (20,1.05) circle (0.15); \fill[funBlue] (20,1.4) circle (0.15);
  \fill[funBlue] (21,0.35) circle (0.15); \fill[funBlue] (21,0.7) circle (0.15); \fill[funBlue] (21,1.05) circle (0.15);
  \fill[funBlue] (22,0.35) circle (0.15); \fill[funBlue] (22,0.7) circle (0.15);
  \fill[funBlue] (23,0.35) circle (0.15);

  % Store B
  \node[left,font=\small\bfseries] at (14,5) {Store B};
  \draw[thick,->] (14.5,3.5) -- (26.5,3.5);
  \foreach \x in {15,16,...,26} { \draw (\x,3.6) -- (\x,3.4); }
  % Store B: 15(1),17(1),18(2),20(2),22(2),24(2),25(1),26(1) = 12
  \fill[funOrange] (15,3.85) circle (0.15);
  \fill[funOrange] (17,3.85) circle (0.15);
  \fill[funOrange] (18,3.85) circle (0.15); \fill[funOrange] (18,4.2) circle (0.15);
  \fill[funOrange] (20,3.85) circle (0.15); \fill[funOrange] (20,4.2) circle (0.15);
  \fill[funOrange] (22,3.85) circle (0.15); \fill[funOrange] (22,4.2) circle (0.15);
  \fill[funOrange] (24,3.85) circle (0.15); \fill[funOrange] (24,4.2) circle (0.15);
  \fill[funOrange] (25,3.85) circle (0.15);
  \fill[funOrange] (26,3.85) circle (0.15);
\end{tikzpicture}
\end{center}

Find the range for each store and explain which store has more consistent hours.
\end{questionText}
\correctAnswer{Store A range $= 5$; Store B range $= 11$. Store A is more consistent.}
\explanation{Store A: $23 - 18 = 5$ hours. Store B: $26 - 15 = 11$ hours. Store A's hours are more tightly clustered, so it is more consistent.}
\end{practiceQuestion}

\begin{practiceQuestion}{m-8-3-q30}{gsa}
\begin{questionText}
The table shows how two classes spend their free time (in percent).

\begin{center}
\begin{tabular}{|l|c|c|}
\hline
\textbf{Activity} & \textbf{Class A} & \textbf{Class B} \\
\hline
Sports & $35\%$ & $20\%$ \\
\hline
Video Games & $25\%$ & $40\%$ \\
\hline
Reading & $15\%$ & $10\%$ \\
\hline
Art/Music & $15\%$ & $20\%$ \\
\hline
Other & $10\%$ & $10\%$ \\
\hline
\end{tabular}
\end{center}

Compare how the two classes spend their free time. Identify the biggest difference and the biggest similarity.
\end{questionText}
\correctAnswer{Biggest difference: Video Games (Class B is $15$ percentage points higher). Biggest similarity: Other ($10\%$ each).}
\explanation{Video Games: $40\% - 25\% = 15$ pp difference (largest). Sports: $35\% - 20\% = 15$ pp (also large). Other: both $10\%$ (identical). Class A favors sports more; Class B favors video games more.}
\end{practiceQuestion}
